\section{Summary}
We demonstrated that the UV divergences of quasi-gluon operators are actually similar to the UV divergences of quasi-quark operators if all divergences are regularized in a gauge invariant scheme.  We proved that UV divergences of all 36 pure quasi-gluon operators are localized in spacetime, and could be multiplicatively renormalized without mixing with each others,
\begin{align}
{\cal O}_{g}^{\mu \nu \rho\sigma}(\xi)=e^{-C_g |\xi_z|}Z_{wg}^{-1}Z_{vg1}^{-s/2}Z_{vg2}^{-(2-s)/2}{\cal O}_{bg}^{\mu \nu \rho\sigma}(\xi),
\end{align}
where $s$ is the number of $z$ components chosen for Lorentz indices $\{ \mu,\nu,\rho,\sigma\}$, and $C_g$, $Z_{wg}$, $Z_{vg1}$ and $Z_{vg2}$ are renormalization constants.
Like the quasi-quark operators, without mixing with each others, hadronic matrix elements of quasi-gluon operators could be examples of good ``lattice cross sections'', as defined in Refs.~\cite{Ma:2014jla,Ma:2017pxb}, which could be calculated in LQCD and factorized into normal PDFs, from which PDFs could be extracted by QCD global analysis of the data of these good ``lattice cross sections'' generated by the first principle LQCD calculations.
